\section{Background \& State of the Art}

\subsection{Host Prediction Problem}
Bacteriophages infect the bacterial hosts through a series of highlt



% The problem of predicting which bacterial hosts a phage can infect from genomic data has been studied computationally 
% for over a decade. Early approaches relied on sequence alignment between phage genomes and bacterial genomes, exploiting 
% the co-evolutionary signal between phage tail proteins and bacterial surface receptors. These methods performed 
% reasonably well when closely related phage-host pairs existed in reference databases, but struggled with novel phages 
% that had no known close relatives.

% More recent approaches framed the problem as a supervised classification task. Given a set of experimentally verified 
% phage-host pairs, a classifier is trained to predict whether a new phage infects a given host based on features derived
% from the phage genome. Among the most informative features are functional annotations of phage proteins: the presence 
% and predicted function of proteins such as lysins, tail fibres, capsid proteins, and portal complexes, since these 
% proteins directly mediate host recognition and cell lysis.

% In this work, the predictions are produced at the protein level. For each protein in a phage genome and each candidate host
% h∈Hh∈H, a classifier outputs a score s∈[0,1]s∈[0,1] reflecting the probability that the protein's function is consistent 
% with infection of that host. Since a phage genome typically contains many proteins — ranging from roughly 10 to 50 in our
% dataset — and each protein may be annotated with one or more functional categories, this yields a structured collection of 
% scores indexed by (protein, host, function) triples. These protein-level scores are then aggregated to the phage level via 
% NaN-aware mean pooling: for each (phage, host, function) combination, we take the mean score across all proteins of that 
% phage annotated with that function, ignoring proteins with missing annotations for that function rather than treating them 
% as zero evidence. The result is a phage-level score vector in [0,1]K[0,1]K, where KK is the number of (host, function) 
% combinations with sufficient data coverage.

% A central empirical observation, discovered during data exploration, is that for the gram-type classification task studied
% in Chapter 5, all score dimensions consistently measure gram-positive evidence: phages known to infect gram-positive 
% bacteria achieve mean scores of approximately 0.810.81 across all KK dimensions, while phages infecting gram-negative 
% bacteria achieve mean scores of approximately 0.080.08. This strong and consistent separation is what makes the gram-type 
% task tractable and serves as a proof-of-concept for the full multi-host framework.




\subsection{The Host Prediction Problem}

% Bacteriophages infect bacterial hosts through a series of highly specific molecular interactions mediated by structural 
% and functional proteins encoded within the phage genome. These interactions determine whether a phage is capable of 
% recognizing, attaching to, and successfully replicating within a particular bacterial strain. The collection of bacterial 
% hosts that a phage can infect is referred to as its \textit{host range}.

% Determining the host range of a phage is one of the central problems in phage biology and has important applications in 
% phage therapy, microbial ecology, and metagenomic analysis. In therapeutic settings, accurate host identification is 
% particularly important because the effectiveness of a phage treatment depends strongly on the compatibility between the 
% selected phage and the target bacterial strain.

% Traditionally, host prediction is performed through experimental laboratory assays in which candidate phages are tested 
% against collections of bacterial hosts. Although these methods are biologically reliable, they are often expensive, 
% time-consuming, and difficult to scale to modern genomic datasets. The rapid increase in sequencing technologies has 
% therefore motivated the development of computational methods capable of predicting phage--host interactions directly from 
% genomic information.

% Computational host prediction remains challenging for several reasons. First, phage genomes are highly diverse and exhibit 
% substantial genetic variability, even among closely related phages. Second, host specificity is rarely determined by a 
% single genomic feature and instead emerges from the combined effect of multiple proteins and functional regions. Finally, 
% genomic datasets frequently contain incomplete annotations, noisy functional predictions, and heterogeneous protein content
% across different phages.

% Modern machine learning approaches therefore increasingly rely on protein-level representations extracted from genomic 
% sequences. In these settings, individual proteins are associated with collections of functional scores or embedding 
% representations that capture biologically relevant information. These protein-level signals are then aggregated to 
% construct phage-level representations suitable for downstream prediction tasks.

% Despite recent progress in computational host prediction, most existing methods continue to produce only deterministic 
% predictions without any associated measure of uncertainty. In biological applications, however, uncertainty quantification 
% is particularly important due to the incomplete and heterogeneous nature of genomic data. This motivates the need for 
% mathematical frameworks capable of producing reliable confidence estimates alongside standard host predictions.
